\editionPage[\pdfbookmark[0]{Démonstrations}{bm-preuves-debut-parent}]{preuves-debut}{Comprendre les preuves}{Identités remarquables et parité}{Les arguments ci-dessous éclairent des propriétés utilisées dans le livre.}
\begin{blocEdition}{preuves-debut}

\begin{demonstration}[Les trois identités remarquables]
Par double distributivité, $(a+b)^2=(a+b)(a+b)=a^2+ab+ba+b^2=a^2+2ab+b^2$.
De même, $(a-b)^2=a^2-ab-ba+b^2=a^2-2ab+b^2$. Enfin
$(a-b)(a+b)=a^2+ab-ab-b^2=a^2-b^2$ : les deux termes opposés s’annulent.
Lues de droite à gauche, ces égalités permettent de factoriser.
\end{demonstration}
\begin{demonstration}[Le carré d’un entier impair est impair]
Un entier impair s’écrit $2n+1$ avec n entier. Son carré vaut
$(2n+1)^2=4n^2+4n+1=2(2n^2+2n)+1$. Le nombre entre parenthèses est entier :
on obtient bien un multiple de 2 augmenté de 1.
\end{demonstration}
\editionRetenir{Une écriture générale prouve ce que quelques essais ne font que suggérer.}

\end{blocEdition}
\editionPage{preuves-2}{Comprendre les preuves}{Pythagore par les aires}{Les arguments ci-dessous éclairent des propriétés utilisées dans le livre.}
\begin{blocEdition}{preuves-2}

Considérons quatre copies d’un triangle rectangle de côtés perpendiculaires
a et b et d’hypoténuse c, disposées dans un carré de côté $a+b$.
\begin{center}\begin{tikzpicture}[scale=.65,font=\small]
\draw[geoTrait](0,0)rectangle(5,5);
\fill[otGreenTint](2,0)--(5,2)--(3,5)--(0,3)--cycle;
\draw[geoTrait](2,0)--(5,2)--(3,5)--(0,3)--cycle;
\node[below]at(1,0){$a$};\node[below]at(3.5,0){$b$};
\node[right]at(5,1){$a$};\node[right]at(5,3.5){$b$};
\node at(2.5,2.5){aire $c^2$};
\end{tikzpicture}\end{center}
Chaque côté du quadrilatère central est une hypoténuse de longueur c.
Ses angles sont droits : les deux angles aigus voisins des triangles
ont une somme de 90 degrés, et complètent un angle plat avec l’angle central.
Le quadrilatère est donc un carré.\par\medskip
L’aire du grand carré vaut $(a+b)^2$. Elle est aussi la somme de l’aire
du carré central et de celles des quatre triangles :
\[
(a+b)^2=c^2+4\times\frac{ab}{2}=c^2+2ab.
\]
En développant puis en soustrayant $2ab$, on obtient
$\res{a^2+b^2=c^2}$.

\end{blocEdition}
\editionPage{preuves-3}{Comprendre les preuves}{Thalès par les aires}{Les arguments ci-dessous éclairent des propriétés utilisées dans le livre.}
\begin{blocEdition}{preuves-3}

Dans le triangle ABC, M appartient à [AB], N à [AC] et $(MN)\parallel(BC)$.
On note $\mathcal A$ l’aire d’un triangle.
\begin{center}\begin{tikzpicture}[scale=.5,font=\small]
\draw[geoTrait](0,0)node[below left]{$A$}--(10,0)node[below right]{$B$}--(4,6)node[above]{$C$}--cycle;
\draw[geoTrait](6,0)node[below]{$M$}--(2.4,3.6)node[left]{$N$};
\draw[otMuted,dashed](6,0)--(4,6) (2.4,3.6)--(10,0);
\end{tikzpicture}\end{center}
\begin{demonstration}[Deux triangles de même aire]
B et C sont sur une droite parallèle à (MN) : ils sont à la même distance
de (MN). Les triangles MNB et MNC ont donc la même base [MN] et la même
hauteur : $\mathcal A(MNB)=\mathcal A(MNC)$. En ajoutant l’aire de AMN,
on obtient $\mathcal A(ANB)=\mathcal A(AMC)$.
\end{demonstration}
\begin{demonstration}[L’égalité des rapports]
Les triangles AMN et ABN ont la même hauteur issue de N ; leurs bases
[AM] et [AB] sont sur la même droite. Donc
$\dfrac{\mathcal A(AMN)}{\mathcal A(ABN)}=\dfrac{AM}{AB}$. De même, avec la
hauteur issue de M, $\dfrac{\mathcal A(AMN)}{\mathcal A(AMC)}=\dfrac{AN}{AC}$.
Les dénominateurs sont égaux d’après la première étape :
$\res{\dfrac{AM}{AB}=\dfrac{AN}{AC}}$.
\end{demonstration}

\end{blocEdition}
\editionPage{preuves-4}{Comprendre les preuves}{Trigonométrie et vecteurs}{Les arguments ci-dessous éclairent des propriétés utilisées dans le livre.}
\begin{blocEdition}{preuves-4}

\begin{demonstration}[$(\cos\widehat B)^2+(\sin\widehat B)^2=1$]
Dans ABC rectangle en A, $\cos\widehat B=AB/BC$ et $\sin\widehat B=AC/BC$. Donc
$(\cos\widehat B)^2+(\sin\widehat B)^2=\dfrac{AB^2+AC^2}{BC^2}$. D’après le théorème
de Pythagore, $AB^2+AC^2=BC^2$ : la somme vaut 1. Connaissant le cosinus
d’un angle aigu, on en déduit son sinus, qui est positif.
\end{demonstration}
\begin{demonstration}[La règle du parallélogramme]
Si ABCD est un parallélogramme, ses côtés opposés [AD] et [BC] sont
parallèles, de même longueur, et parcourus dans le même sens :
$\vect{AD}=\vect{BC}$. Par la relation de Chasles,
$\vect{AB}+\vect{AD}=\vect{AB}+\vect{BC}=\vect{AC}$.
\end{demonstration}
\editionRetenir{Une égalité de vecteurs traduit à la fois un parallélisme, un sens et une égalité de longueurs.}

\end{blocEdition}
